\documentclass[10pt,a4paper]{amsart}
\usepackage{amsmath,amssymb,amsthm,graphicx}
% Retained minimal preamble for the Stacks Project chapter more-algebra.tex (More on Algebra).
% The chapter file pinned-original-source.tex carries only \input{preamble} on its first line, so no
% preamble text lives inside the chapter itself.  The Stacks Project uses one single shared preamble for
% every chapter; the snapshot stages it identically as sources/stacks-chow/preamble.tex and
% sources/stacks-etale-cohomology/preamble.tex.  The block below is copied mechanically, line for line,
% from sources/stacks-chow/preamble.tex lines 14-18 (the Xy-pic setup) and 154-217 (the native theorem
% declarations together with the mathematical macros).  Nothing was added, removed or redefined; the
% wrapper class and the amsmath/amssymb/amsthm/graphicx base packages are supplied by the assembler.
% Chapter: more-algebra.tex -- More on Algebra.
% For commutative diagrams we use Xy-pic
\usepackage[all]{xy}

% We use 2cell for 2-commutative diagrams.
\xyoption{2cell}

% Theorem environments.
%
\theoremstyle{plain}
\newtheorem{theorem}[subsection]{Theorem}
\newtheorem{proposition}[subsection]{Proposition}
\newtheorem{lemma}[subsection]{Lemma}

\theoremstyle{definition}
\newtheorem{definition}[subsection]{Definition}
\newtheorem{example}[subsection]{Example}
\newtheorem{exercise}[subsection]{Exercise}
\newtheorem{situation}[subsection]{Situation}

\theoremstyle{remark}
\newtheorem{remark}[subsection]{Remark}
\newtheorem{remarks}[subsection]{Remarks}

\numberwithin{equation}{subsection}

% Macros
%
\def\lim{\mathop{\mathrm{lim}}\nolimits}
\def\colim{\mathop{\mathrm{colim}}\nolimits}
\def\Spec{\mathop{\mathrm{Spec}}}
\def\Hom{\mathop{\mathrm{Hom}}\nolimits}
\def\Ext{\mathop{\mathrm{Ext}}\nolimits}
\def\SheafHom{\mathop{\mathcal{H}\!\mathit{om}}\nolimits}
\def\SheafExt{\mathop{\mathcal{E}\!\mathit{xt}}\nolimits}
\def\Sch{\mathit{Sch}}
\def\Mor{\mathop{\mathrm{Mor}}\nolimits}
\def\Ob{\mathop{\mathrm{Ob}}\nolimits}
\def\Sh{\mathop{\mathit{Sh}}\nolimits}
\def\NL{\mathop{N\!L}\nolimits}
\def\CH{\mathop{\mathrm{CH}}\nolimits}
\def\proetale{{pro\text{-}\acute{e}tale}}
\def\etale{{\acute{e}tale}}
\def\QCoh{\mathit{QCoh}}
\def\Ker{\mathop{\mathrm{Ker}}}
\def\Im{\mathop{\mathrm{Im}}}
\def\Coker{\mathop{\mathrm{Coker}}}
\def\Coim{\mathop{\mathrm{Coim}}}

% Boxtimes
%
\DeclareMathSymbol{\boxtimes}{\mathbin}{AMSa}{"02}

%
% Macros for moduli stacks/spaces
%
\def\QCohstack{\mathcal{QC}\!\mathit{oh}}
\def\Cohstack{\mathcal{C}\!\mathit{oh}}
\def\Spacesstack{\mathcal{S}\!\mathit{paces}}
\def\Quotfunctor{\mathrm{Quot}}
\def\Hilbfunctor{\mathrm{Hilb}}
\def\Curvesstack{\mathcal{C}\!\mathit{urves}}
\def\Polarizedstack{\mathcal{P}\!\mathit{olarized}}
\def\Complexesstack{\mathcal{C}\!\mathit{omplexes}}
% \Pic is the operator that assigns to X its picard group, usage \Pic(X)
% \Picardstack_{X/B} denotes the Picard stack of X over B
% \Picardfunctor_{X/B} denotes the Picard functor of X over B
\def\Pic{\mathop{\mathrm{Pic}}\nolimits}
\def\Picardstack{\mathcal{P}\!\mathit{ic}}
\def\Picardfunctor{\mathrm{Pic}}
\def\Deformationcategory{\mathcal{D}\!\mathit{ef}}
\title{Finite presentation relative to R survives localization: M of finite presentation over R\_f becomes M\_g of finite presentation over R}
\author{The Stacks Project Authors}
\date{stacks-project:more-algebra.tex:local-snapshot-20261012; GNU FDL 1.2 or later}
\begin{document}
\maketitle
\setcounter{subsection}{1}
\noindent\textit{Editorial note.} All delivered mathematical text is the author's own, copied from the pinned Stacks Project chapter file more-algebra.tex; no statement or proof step was rewritten, reordered or removed. The counted statement is the labeled lemma environment at source lines 21014-21020 (canonical tag 065A) and the counted proof is the proof environment at source lines 21022-21045, which immediately follows it and is its own proof.\par\medskip
\noindent\textit{Editorial note.} The last sentence of the proof as printed says that the displayed map is a presentation of M\_f, while the construction is set up to present M\_g. That slip is copied verbatim and is not corrected, since the lane delivers the author's own text and does not edit mathematics.\par\medskip
\noindent\textit{Editorial note.} The chapter file pinned-original-source.tex carries only an input of the shared preamble on its first line, so no preamble text lives inside the chapter. The retained file sources/stacks-more-algebra/retained-preamble.tex is a mechanical, line-for-line copy of the project's single shared preamble as staged in the snapshot (Xy-pic setup plus the native theorem declarations and mathematical macros). Nothing was added, removed or redefined, and the excerpt serves only as preamble, never as a span coordinate.\par\medskip
\noindent\textit{Editorial note.} The source URL is the source-level Stacks Project URL, as required when more than one card is produced from one source; the per-claim tag is carried by claim\_key and by this note. The counted statement and proof refer to no labelled item of the chapter and cite no bibliography entry, so this card contributes no external reference stub and needs no bibliography.\par\medskip
\noindent\textit{Allowed standard prerequisites (editorial).}\begin{itemize}
\item The localization R\_f is the quotient of the polynomial ring in one variable by fx\_0 - 1, so a presentation over R\_f can be rewritten as a presentation over R with one extra variable. used at proof lines 21023-21026, where R\_f = R[x\_0]/(fx\_0 - 1) is recorded and the presentation is extended by the variables x\_0 and x\_\{n+1\}. Produces the ambient polynomial ring over R in which the enlarged presentation of the localized module is written.
\item Elements and matrix entries of a finitely generated algebra over R\_f can be lifted to the polynomial ring over R. used at proof lines 21027-21034, where g' and the entries h'\_ij are chosen to map to g and to h\_ij. Supplies the blocks of the lifted matrix, which is what makes the presentation a presentation over R rather than over R\_f.
\item A surjection of free modules with a given matrix, displayed with three blocks, presents the localized module once the second block encodes fx\_0 - 1 and the third block encodes x\_\{n+1\}g' - 1. used at proof lines 21037-21044, where the t by (s + 2t) matrix is described block by block. Verifies that the displayed map is onto M\_g, which is the conclusion of the lemma.
\end{itemize}\medskip
\begin{lemma}
\label{lemma-localize-relative-finite-presentation}
Let $R$ be a ring, $f \in R$ an element, $R_f \to A$ is a finite type ring map,
$g \in A$, and $M$ an $A$-module. If $M$ of finite presentation relative
to $R_f$, then $M_g$ is an $A_g$-module of finite presentation relative
to $R$.
\end{lemma}

\begin{proof}
Choose a presentation $R_f[x_1, \ldots, x_n] \to A$. We write
$R_f = R[x_0]/(fx_0 - 1)$. Consider the presentation
$R[x_0, x_1, \ldots, x_n, x_{n + 1}] \to A_g$ which extends the given
map, maps $x_0$ to the image of $1/f$, and maps $x_{n + 1}$ to $1/g$.
Choose $g' \in R[x_0, x_1, \ldots, x_n]$ which maps to $g$ (this is
possible). Suppose that
$$
R_f[x_1, \ldots, x_n]^{\oplus s} \to
R_f[x_1, \ldots, x_n]^{\oplus t} \to M \to 0
$$
is a presentation of $M$ given by a matrix $(h_{ij})$. Pick
$h'_{ij} \in R[x_0, x_1, \ldots, x_n]$ which map to $h_{ij}$.
Then
$$
R[x_0, x_1, \ldots, x_n, x_{n + 1}]^{\oplus s + 2t} \to
R[x_0, x_1, \ldots, x_n, x_{n + 1}]^{\oplus t} \to M_g \to 0
$$
is a presentation of $M_f$.
Here the $t \times (s + 2t)$ matrix defining the map has a first
$t \times s$ block consisting of the matrix $h'_{ij}$, a second
$t \times t$ block which is $(x_0f - )I_t$, and a third block
which is $(x_{n + 1}g' - 1)I_t$.
\end{proof}

\end{document}
